% The Linnik exponent paper: a skeleton. Each section says what it will contain; the claims table of
% Section 7 is complete and names the Lean theorem behind each claim. The paper claims only its own steps: the
% results it uses from the literature are documented in formalisations/ and cited, not restated as its own.
\documentclass[11pt]{article}
\usepackage[T1]{fontenc}
\usepackage[utf8]{inputenc}
\usepackage{amsmath,amssymb,amsthm}
\usepackage{booktabs}
\usepackage[hidelinks]{hyperref}

\newtheorem{theorem}{Theorem}[section]
\newtheorem{proposition}[theorem]{Proposition}
\newtheorem{lemma}[theorem]{Lemma}
\theoremstyle{definition}
\newtheorem{definition}[theorem]{Definition}
\newtheorem{hypothesis}{Hypothesis}
\renewcommand{\thehypothesis}{H\arabic{hypothesis}}

\newcommand{\lean}[1]{\texttt{#1}}

\title{The least prime in an arithmetic progression:\\ the exponent $899/400$}
% The authors are to be filled in.
\author{}
\date{Draft}

\begin{document}
\maketitle

\begin{abstract}
We prove that there is a constant $K$ such that, for every $q \ge 1$ and every $a$ prime to $q$, there is a prime
$p \equiv a \pmod q$ with $p \le K q^{899/400}$. The proof is a sieve argument. Its Type~II estimate uses two
recent results: that no Dirichlet $L$-function vanishes in $\operatorname{Re} s > 11/12$, from OpenAI's release, and
Chen, Gupta and Li's zero-density estimate. The theorem, these two results and the classical results the proof uses
are proved in Lean~4, and the proof depends on no axioms beyond Lean's standard three.
\end{abstract}

\section{Introduction}\label{sec:intro}

\begin{theorem}\label{thm:main}
There is $K$ such that for every $q \ge 1$ and every $a$ prime to $q$ there is a prime $p \equiv a \pmod q$ with
$p \le K q^{899/400}$.
\end{theorem}

Linnik \cite{linnik44} proved that there are absolute constants $K$ and $L$ such that every reduced residue class
mod $q$ contains a prime $p \le K q^L$. The admissible exponent $L$ has been lowered many times. Heath-Brown
\cite{heath-brown92} reached $L = 5.5$ and Xylouris $L = 5$ \cite{xylouris11}, and then $L < 5$ \cite{xylouris18}.
Under the generalised Riemann hypothesis every $L > 2$ is admissible \cite{lls}. Theorem~\ref{thm:main} is
unconditional, with $L = 899/400 = 2.2475$.

\subsection*{The zeros near $\operatorname{Re} s = 1$}

The difficulty in Linnik's theorem is the zeros of the $L$-functions mod $q$ near the line $\operatorname{Re} s = 1$.
By the explicit formula, a zero $\rho = \beta + i\gamma$ of $L(s, \chi)$ contributes about $x^{\beta - 1}$, relative to
the main term, to the count of primes $p \le x$ in a class mod $q$. A zero-density estimate with exponent $A$, a
bound of about $(qT)^{A(1 - \sigma)}$ for the number of zeros with real part at least $\sigma$, makes the zeros of
real part about $\sigma$ contribute about $(q^A/x)^{1 - \sigma}$ together. Classically a zero may have
$1 - \beta$ as small as $c/\log q$, and at $x = q^L$ the contribution $(q^A/x)^{1 - \beta}$ is then a constant, not a
small quantity. The record proofs control these zeros with zero-free regions, density estimates without logarithmic
factors and the Deuring--Heilbronn phenomenon, and the constants of all three enter the exponent.

The zero-free half-plane of OpenAI's release \cite{openai}, Hypothesis~\ref{hyp:half-plane}, removes the zeros with
$\beta > 11/12$. Every remaining zero then contributes at most a power $x^{-1/12}$, and the problem becomes one of
density: a density exponent $A$ for $1/2 \le \sigma \le 11/12$ gives Linnik's theorem with every exponent $L > A$.
Huxley's density estimate \cite{huxley75} gives $A = 12/5$. Chen, Gupta and Li \cite{cgl} prove $A = 7/3$, and
their Corollary~1.4, on primes in short intervals in arithmetic progressions under a zero-free region that
Hypothesis~\ref{hyp:half-plane} supplies, gives the exponent $7/3 + \varepsilon$.
Section~\ref{sec:comparison} gives the details.

\subsection*{Below the density exponent}

This paper goes below $7/3$ with a sieve. Harman's sieve \cite{harman83, harman} bounds from below the number of
primes in the set $\mathcal A$ of integers $n \le x$ with $n \equiv a \pmod q$ by comparing it with the set
$\mathcal B$ of all integers $n \le x$, weighted by $1/\varphi(q)$. It needs two kinds of information. Type~I
information is the comparison of $\mathcal A$ with $\mathcal B$ on the multiples of each small integer, which
orthogonality gives. Type~II information is the comparison on bilinear sums
$\sum_{mn} a_m b_n$, with $m$ in a range $(x^{u_1}, x^{u_2})$ and arbitrary bounded coefficients, which needs
the zeros. A comparison over all such sums is not available at $x = q^{899/400}$. What is available is a window
$[u_1, u_2]$ of exponents on which it holds, and the sieve recovers the primes from the sums it can compare,
losing a fraction of them that a finite computation bounds.

The paper's contributions are three.
\begin{enumerate}
\item \emph{A Type~II estimate over the zeros' real parts} (Section~\ref{sec:type2}). The characters mod $q$ are
  sorted into bins by the largest real part of their $L$-function's zeros near a given height. In each bin the
  bilinear sum is bounded by the zeros, which make the Dirichlet polynomial of the primes small, and by the number
  of characters in the bin, which a density estimate bounds; the best of orthogonality and the
  Hal\'asz--Montgomery inequality is taken in each bin. Hypotheses~\ref{hyp:half-plane} and~\ref{hyp:density},
  with the classical density estimates of Ingham and Huxley, give the window
  $[u_1, u_2] = [41335/100000,\ 48069/100000]$ at $L = 899/400$.
\item \emph{A decomposition by Buchstab's identity, with a refinement} (Section~\ref{sec:certificate}). The
  identity is applied to depth four. The terms in which the product of some of the prime variables has its
  exponent in the window, or in its reflection $[1 - u_2, 1 - u_1]$, are Type~II sums and are evaluated, and the
  others are dropped. The refinement keeps, in a dropped term of depth two, the part whose integers have exactly two
  prime factors in the cofactor and are, in some grouping of their four prime factors, a Type~II sum.
\item \emph{A certificate} (Section~\ref{sec:certificate}). The fraction of the primes the decomposition loses is
  bounded above by a sum over the cells of a grid of the exponents, computed exactly in integers. On a grid of
  $300$ cells a side the sum is below $0.99386$, which is less than $1$, so a positive proportion of the primes
  remains.
\end{enumerate}

The rest is not new. Hypotheses~\ref{hyp:half-plane} and~\ref{hyp:density} are OpenAI's and Chen, Gupta and Li's,
and the density estimates of Ingham \cite{ingham} and Huxley \cite{huxley72, huxley75}, the Hal\'asz--Montgomery
inequality \cite{halasz, montgomery69} and Harman's sieve are classical. Section~\ref{sec:inputs} states the inputs
and Section~\ref{sec:type1} the form of the sieve that is used.

\subsection*{The formalisation}

Every step of the proof, and every result it uses from the literature, is proved in Lean~4 with the Mathlib
library. The Lean theorem \lean{Linnik22475R.linnik\_22475r} is Theorem~\ref{thm:main}, and its axioms are Lean's
standard three. The certificate's sum is checked by Lean's kernel in 534 parts. Section~\ref{sec:lean} names the Lean
theorem behind each claim of the paper, and the directory \texttt{formalisations/} of the repository documents each
result from the literature: its source, its Lean statement and the differences between them.

\subsection*{Outline}

Section~\ref{sec:inputs} states the two recent inputs and the density estimates the Type~II estimate uses.
Section~\ref{sec:type2} proves the Type~II estimate, and Section~\ref{sec:type1} gives the Type~I estimate and the
sieve. Section~\ref{sec:certificate} describes the decomposition and its certificate, and Section~\ref{sec:proof}
assembles the proof. Section~\ref{sec:lean} describes the formalisation, and Section~\ref{sec:comparison} compares
Theorem~\ref{thm:main} with what the inputs give without the sieve.

\section{Notation and the two inputs}\label{sec:inputs}

Throughout, $q \ge 1$ is the modulus and $a$ is prime to $q$. The Dirichlet characters mod $q$ are $\chi$, the trivial
one is $\chi_0$, and $L(s, \chi)$ is the $L$-function of $\chi$, continued to the whole plane (in Lean, Mathlib's
\lean{DirichletCharacter.LFunction}). For $1/2 \le \sigma \le 1$ and $T \ge 2$ let
\[
  N^*(\sigma, T; q) = \#\{\chi \bmod q : L(s, \chi) = 0 \text{ for some } s \text{ with }
    \operatorname{Re} s \ge \sigma,\ |\operatorname{Im} s| \le T\}
\]
be the number of characters mod $q$ whose $L$-function has a zero in the rectangle. It counts characters, not zeros,
so a bound for the number of zeros bounds it. The proof is a sieve on the integers up to $x = q^L$, with $L = 899/400$
in the end, and $q$ is taken large.

\begin{hypothesis}[the zero-free half-plane]\label{hyp:half-plane}
For every $q \ge 1$, every Dirichlet character $\chi$ mod $q$ and every $s$ with $\operatorname{Re} s > 11/12$,
$L(s, \chi) \ne 0$, except at $s = 1$ for the trivial character.
\end{hypothesis}

This is OpenAI's result \cite{openai}: their companion preprint proves it, and their main preprint proves the
stronger $\operatorname{Re} s > 7/8$. It is proved in Lean as \lean{Eis.LFunction\_ne\_zero\_of\_re\_gt}, following
the companion preprint. \texttt{formalisations/openai-half-plane.md} gives the attribution and the differences
between the Lean statement and the source's.

\begin{hypothesis}[a zero-density estimate]\label{hyp:density}
For every $\varepsilon > 0$ there is $C$ such that for every $q \ge 1$, every $5/7 \le \sigma \le 16/21$ and every
$T \ge 2$, at most $C (qT)^{(7/3)(1 - \sigma) + \varepsilon}$ characters mod $q$ have an $L$-function with a zero in
$\operatorname{Re} s \ge \sigma$, $|\operatorname{Im} s| \le T$.
\end{hypothesis}

This follows from Chen, Gupta and Li's estimate \cite{cgl}. Their estimate counts zeros, which gives the count of
characters stated here. It is proved in Lean as \lean{LinnikCGLDensity.cglDensity}, and \texttt{formalisations/}
documents it.

\subsection*{The density curve}

The Type~II estimate of Section~\ref{sec:type2} needs a bound for $N^*(\sigma, T; q)$ at every $\sigma$ in
$[1/2, 1]$. It takes it in the form of a curve.

\begin{definition}\label{def:curve}
A function $g : \mathbb R \to \mathbb R$ is a \emph{density curve} if there is $C \ge 0$ such that for every
$q \ge 1$, every $1/2 \le \sigma \le 1$ and every $T \ge 2$,
\[
  N^*(\sigma, T; q) \le C\, (qT)^{g(\sigma)} \bigl(1 + \log qT\bigr)^{10}.
\]
\end{definition}

The curve of this paper is the function $g_7$ defined as follows; each of its pieces is a known density estimate,
and the steps make each piece a constant multiple of $1 - \sigma$.
\begin{itemize}
\item $g_7(\sigma) = 1$ for $\sigma < 4/7$.
\item On $[4/7, 5/7]$, divided into twenty steps $[\sigma_1, \sigma_2)$ of length $1/140$,
  $g_7(\sigma) = \bigl(3/(2 - \sigma_2) + 1/2000\bigr)(1 - \sigma)$: Ingham's exponent at the step's right end.
\item On $[5/7, 16/21)$, $g_7(\sigma) = (7/3 + 1/2000)(1 - \sigma)$.
\item On $[16/21, 5/6]$, divided into twenty steps $[\sigma_1, \sigma_2)$ of length $1/280$,
  $g_7(\sigma) = \bigl(3/(3\sigma_1 - 1) + 1/2000\bigr)(1 - \sigma)$: Huxley's exponent at the step's left end.
\item $g_7(\sigma) = (101/50)(1 - \sigma)$ for $\sigma \ge 5/6$.
\end{itemize}
The exponents are written in Lean as fractions; the first step's is $3/(2 - 81/140) + 1/2000 = 840199/398000$. On
$[4/7, 16/21]$ the exponent rises from about $2.111$ to $7/3 + 1/2000$, and on $[16/21, 5/6]$ it falls back to
about $2.015$. On each step $A(1 - \sigma) \le A(1 - \sigma_1) \le 1$, so $g_7 \le 1$ everywhere, and
$g_7(1/2) = 1$.

\begin{proposition}\label{prop:curve}
Under Hypothesis~\ref{hyp:density}, $g_7$ is a density curve.
\end{proposition}

\begin{proof}
A finite union of ranges on each of which the bound holds, with its own constant, is a range on which it holds with
the sum of the constants, so it suffices to prove the bound on each piece.
\begin{itemize}
\item For $\sigma < 4/7$: $N^*(\sigma, T; q) \le \varphi(q) \le qT$.
\item On each step of $[4/7, 5/7]$: Ingham's estimate \cite{ingham} in the form over the characters mod $q$, which
  bounds $N^*(\sigma, T; q)$ by $C(qT)^{A(1 - \sigma)}(\log qT)^{10}$ on $[\sigma_1, \sigma_2]$ for every $A$
  above $3/(2 - \sigma_2)$ (\texttt{formalisations/ingham-density.md}). Ingham proved his estimate for the zeros of
  $\zeta$; the form over the characters is Theorem~12.1 of Montgomery's book \cite{montgomery-topics}.
\item On $[5/7, 16/21]$: Hypothesis~\ref{hyp:density} with $\varepsilon = (1/2000)(5/21)$, since $1 - \sigma \ge 5/21$
  there.
\item On each step of $[16/21, 5/6]$: Huxley's estimate \cite{huxley75} over the characters mod $q$, with the
  exponent $3/(3\sigma_1 - 1) + 1/2000$ on $[\sigma_1, \sigma_2]$ (\texttt{formalisations/huxley-density.md}).
\item For $\sigma \ge 5/6$: the exponent $101/50$, by the same large-values argument at fixed parameters on
  $[5/6, 11/12]$, and by Hypothesis~\ref{hyp:half-plane} above $11/12$, where at most the trivial character is
  counted (\texttt{formalisations/huxley-density.md}).
\end{itemize}
Each of these bounds is a theorem of the library, with a factor $(\log qT)^{10}$ at most $(1 + \log qT)^{10}$; the
claims table of Section~\ref{sec:lean} names them.
\end{proof}

The only piece that is not classical is the one on $[5/7, 16/21]$, where Chen, Gupta and Li's exponent $7/3$ replaces
Huxley's, which is $12/5$ at $\sigma = 3/4$. Hypothesis~\ref{hyp:half-plane} enters the curve only above $11/12$, and
enters the Type~II estimate through the largest real part of the zeros, which it bounds by $11/12$.

\section{The Type II estimate}\label{sec:type2}

% To be written: smooth prime sums bounded through each character's zeros near the height (Lemma 1 of the route);
% the bins of characters by the real parts of their zeros; the top bin by Hypothesis H1; the other bins by the
% density curve; the Halasz-Montgomery branch on the top piece [5/6, 11/12] (formalisations/halasz-montgomery.md,
% cited to Halasz and Montgomery); the window
% [41335/100000, 48069/100000] L at L = 899/400.

\section{Type I estimates and Harman's sieve}\label{sec:type1}

% To be written: the Type I estimate by orthogonality, the fundamental lemma, and Harman's comparison principle in
% the form used, cited to Harman, Brun and Buchstab (formalisations/harman-sieve.md) and not restated as new.

\section{The decomposition and its certificate}\label{sec:certificate}

% To be written: Buchstab's identity on the cells of the grid (n = 300; the window's ends rounded inward over
% DEN = 2^40); the kept cells and the discarded ones; the loss as a sum over the cells, bounded above with the
% scale K = 2^60; the refined certificate and the kernel's check of its total.

\section{Proof of Theorem~\ref{thm:main}}\label{sec:proof}

% To be written: the route theorem (Section 7's table) at the parameters, with Hypotheses H1 and H2.

\section{The formalisation}\label{sec:lean}

Every claim of the paper is a theorem of the Lean library in \texttt{lib/}. The table names each one. The file
\texttt{lib/check/Check.lean} prints the axioms of these theorems, and \texttt{scripts/check\_axioms.py} checks that
they are \lean{propext}, \lean{Classical.choice} and \lean{Quot.sound} only.

\begin{center}
\begin{tabular}{@{}ll@{}}
\toprule
Claim & Lean \\
\midrule
Theorem~\ref{thm:main} & \lean{Linnik22475R.linnik\_22475r} \\
Theorem~\ref{thm:main} under Hypothesis~\ref{hyp:density} & \lean{Linnik22475R.linnik\_22475r\_of} \\
Proposition~\ref{prop:curve}: $g_7$ is a density curve & \lean{SieveTypeIIU226C.counts\_7c} \\
The route theorem (Section~\ref{sec:proof}) & \lean{SieveMain2.linnik\_route\_lt2} \\
The Type~II estimate on the window (Section~\ref{sec:type2}) & \lean{SieveTypeIIU22475C.typeII\_22475c} \\
The certificate's total is below $K$ (Section~\ref{sec:certificate}) & \lean{SieveCert.total2\_lt\_22475R} \\
\midrule
Hypothesis~\ref{hyp:half-plane} (OpenAI; \texttt{formalisations/}) & \lean{Eis.LFunction\_ne\_zero\_of\_re\_gt} \\
Hypothesis~\ref{hyp:density} (Chen--Gupta--Li; \texttt{formalisations/}) & \lean{LinnikCGLDensity.cglDensity} \\
Ingham's density estimate (\texttt{formalisations/}) & \lean{LinnikInghamC.charDensity\_ingham\_c} \\
Huxley's density estimate (\texttt{formalisations/}) & \lean{LinnikHuxStep.charDensity\_gen} \\
Huxley's estimate from $5/6$ on (\texttt{formalisations/}) & \lean{LinnikStep.charDensity\_step} \\
The Hal\'asz--Montgomery inequality (\texttt{formalisations/}) & \lean{LinnikPower.hm\_duality} \\
Harman's sieve in the form used (\texttt{formalisations/}) & \lean{SieveProp6.prop6} \\
\bottomrule
\end{tabular}
\end{center}

% To be written: how faithful the Lean statements are (Mathlib's DirichletCharacter.LFunction, the counts of
% characters), and how to build and check (the repository's README).

\section{What the inputs give without the sieve}\label{sec:comparison}

% To be written: Hypotheses H1 and H2 by the classical deduction give 7/3 + epsilon (Chen-Gupta-Li's Corollary 1.4);
% the library's own deduction (Linnik.linnik_of_charDensity) and the exponents it gives; the trivial count gives
% 12 + epsilon from H1 alone. This section claims nothing new; it places Theorem 1.1.

% The entries are to be checked against the sources before submission.
\begin{thebibliography}{99}
\bibitem{linnik44} U. V. Linnik, \emph{On the least prime in an arithmetic progression. I. The basic theorem},
  Rec. Math. [Mat. Sbornik] N.S. 15(57) (1944), 139--178.
\bibitem{heath-brown92} D. R. Heath-Brown, \emph{Zero-free regions for Dirichlet $L$-functions, and the least prime
  in an arithmetic progression}, Proc. London Math. Soc. (3) 64 (1992), 265--338.
\bibitem{xylouris11} T. Xylouris, \emph{\"Uber die Nullstellen der Dirichletschen $L$-Funktionen und die kleinste
  Primzahl in einer arithmetischen Progression}, Bonner Math. Schriften 404, Bonn, 2011.
\bibitem{xylouris18} T. Xylouris, \emph{Linniks Konstante ist kleiner als 5}, Chebyshevskii Sb. 19 (2018), 80--94.
\bibitem{lls} Y. Lamzouri, X. Li and K. Soundararajan, \emph{Conditional bounds for the least quadratic non-residue
  and related problems}, Math. Comp. 84 (2015), 2391--2412.
\bibitem{openai} OpenAI, \emph{The Quasi-Riemann Hypothesis: a zero-free half-plane $\operatorname{Re}(s) > 7/8$},
  preprint, 30 September 2026, and its companion preprint of 5 October 2026; \url{https://github.com/openai/math},
  entry 003.
\bibitem{cgl} B. Chen, V. Gupta and Y. C. Li, \emph{Large value estimates for Dirichlet polynomials with characters
  and zero density of Dirichlet $L$-functions}, arXiv:2507.08296v2.
\bibitem{ingham} A. E. Ingham, \emph{On the estimation of $N(\sigma, T)$}, Quart. J. Math. Oxford Ser. 11 (1940),
  201--202.
\bibitem{montgomery-topics} H. L. Montgomery, \emph{Topics in Multiplicative Number Theory}, Lecture Notes in Math.
  227, Springer, Berlin, 1971.
\bibitem{huxley72} M. N. Huxley, \emph{On the difference between consecutive primes}, Invent. Math. 15 (1972),
  164--170.
\bibitem{huxley75} M. N. Huxley, \emph{Large values of Dirichlet polynomials, III}, Acta Arith. 26 (1975), 435--444.
\bibitem{halasz} G. Hal\'asz, \emph{\"Uber die Mittelwerte multiplikativer zahlentheoretischer Funktionen}, Acta Math.
  Acad. Sci. Hungar. 19 (1968), 365--403.
\bibitem{montgomery69} H. L. Montgomery, \emph{Mean and large values of Dirichlet polynomials}, Invent. Math. 8
  (1969), 334--345.
\bibitem{harman83} G. Harman, \emph{On the distribution of $\alpha p$ modulo one}, J. London Math. Soc. (2) 27 (1983),
  9--18.
\bibitem{harman} G. Harman, \emph{Prime-Detecting Sieves}, London Math. Soc. Monographs 33, Princeton University
  Press, 2007.
\bibitem{brun} V. Brun, \emph{Le crible d'Eratosth\`ene et le th\'eor\`eme de Goldbach}, Videnskapsselskapets
  Skrifter I, Mat.-naturv. Klasse, Kristiania, 1920, no.~3.
\bibitem{buchstab} A. A. Buchstab, \emph{Asymptotische Absch\"atzung einer allgemeinen zahlentheoretischen
  Funktion}, Mat. Sb. 2(44) (1937), 1239--1246.
\end{thebibliography}

\end{document}
